\subsection{取值于李代数的微分形式的原函数}

引理 5. 设 X 是一个 \(\mathbf{C}^r\) 类流形，F 和 \(\mathbf{F}'\) 是以 X 为底空间的 \(\mathbf{C}^r\) 类向量丛，且 \(\phi\) 是从 F 到 \(\mathbf{F}'\) 的态射。对每个 \(x \in \mathbf{X}\)，令 \(\mathbf{S}_x\) 为由下列元素组成的集合：

\[
(a, \phi (a)) \in \mathrm{F} _ {x} \oplus \mathrm{F} _ {x} ^ {\prime}
\]

其中 \(a \in \mathbf{F}_x\)。则 \(\mathbf{S}\) 是各 \(\mathbf{S}_x\) 的并，且是 \(\mathbf{F} \oplus \mathbf{F}'\) 的一个向量子丛。

令 \(\theta\) 和 \(\theta'\) 为从 \(F \oplus F'\) 到自身的映射，定义如下：若 \((u, v) \in F_x \oplus F_x'\)，则

\[
\theta (u, v) = (u, v + \phi (u)), \quad \theta^ {\prime} (u, v) = (u, v - \phi (u)).
\]

由 Differentiable and Analytic Manifolds, R, 7.7.1，\(\theta\) 和 \(\theta'\) 是从 \(F \oplus F'\) 到自身的态射。显然 \(\theta \circ \theta' = \theta' \circ \theta = \mathrm{Id}_{F\oplus F'}\)。因此 \(\theta\) 和 \(\theta'\) 是 \(F \oplus F'\) 的自同构。于是，\(S = \theta(F \oplus \{0\})\) 是 \(F \oplus F'\) 的一个向量子丛。

引理 6. 设 G 是一个李群芽，\(\omega\) 是 G 的左规范微分形式（§ 3, no. 18.9），M 是一个 \(\mathbf{C}^r\) 类流形（\(r \geqslant 2\)），而 \(\alpha\) 是 M 上取值于 L(G) 的 \(\mathbf{C}^{r-1}\) 类 1 次微分形式。

\begin{enumerate}
\def\labelenumi{(\roman{enumi})}
\tightlist
\item
  \(\mathbf{T}(\mathbf{M}\times \mathbf{G})\) 中使微分形式
\end{enumerate}

\[
\theta = \mathrm{pr} _ {1} ^ {*} \alpha - \mathrm{pr} _ {2} ^ {*} \omega
\]

为零的元素构成 T(M × G) 的一个 \(\mathbf{C}^{r-1}\) 类向量子丛 S。

\begin{enumerate}
\def\labelenumi{(\roman{enumi})}
\setcounter{enumi}{1}
\item
  对所有 \((x,g)\in \mathbf{M}\times \mathbf{G},\mathrm{T}(\mathrm{pr}_1)|\mathrm{S}_{(x,g)}\) 都是从 \(\mathrm{S}_{(x,g)}\) 到 \(\mathrm{T}_x(\mathbf{M})\) 的同构。
\item
  若 \(d\alpha + [\alpha]^2 = 0\)（cf.~§3, no. 14），则向量子丛 S 可积。
\end{enumerate}

若 \((x,g)\in \mathbf{M}\times \mathbf{G}\) 且 \((u,v)\in \mathrm{T}_x(\mathrm{M})\times \mathrm{T}_g(\mathrm{G})\)，则

\[
\theta_ {(x, g)} (u, v) = \alpha (u) - g ^ {- 1} v.
\]

因此，\(\theta_{(x,g)}\) 的核是集合 \(\mathrm{S}_{(x,g)}\)，它由 \((u,g\alpha (u))\)（其中 \(u\in \mathrm{T}_x(\mathbf{M})\)）组成，从而得到 (ii)。我们把 \(\mathrm{T}(\mathrm{M}\times \mathrm{G})\) 看作两个向量丛 F 和 \(\mathbf{F}'\) 的直和，其中 \(\mathrm{F}_{(x,g)} = \mathrm{T}_x(\mathrm{M})\times \{0\}\) 且 \(\mathrm{F}_{(x,g)}^{\prime} = \{0\} \times \mathrm{T}_g(\mathrm{G})\)，对所有

\[
(x, g) \in \mathbf {M} \times \mathbf {G}.
\]

对 \(u \in \mathrm{T}_{x}(\mathbf{M}) \times \{0\}\)，记 \(\phi(u) = (0, g\alpha(u))\)。利用 \(\mathrm{T}(\mathrm{G})\) 的左平凡化可见，\(\phi\) 是从 \(\mathbf{F}\) 到 \(\mathbf{F}'\) 的态射，从而得到 (i)（引理 5）。最后，若 \(d\alpha + [\alpha]^2 = 0\)，则

\[
\begin{array}{r l} {d \theta = \mathrm{pr} _ {1} ^ {*} (d \alpha) - \mathrm{pr} _ {2} ^ {*} (d \omega)} \\ & {= - \frac {1}{2} (\mathrm{pr} _ {1} ^ {*} \alpha \wedge \mathrm{pr} _ {1} ^ {*} \alpha - \mathrm{pr} _ {2} ^ {*} \omega \wedge \mathrm{pr} _ {2} ^ {*} \omega)} \\ & {= - \frac {1}{2} (\mathrm{pr} _ {1} ^ {*} \alpha - \mathrm{pr} _ {2} ^ {*} \omega) \wedge (\mathrm{pr} _ {1} ^ {*} \alpha + \mathrm{pr} _ {2} ^ {*} \omega)} \\ & {= - \frac {1}{2} \theta \wedge (\mathrm{pr} _ {1} ^ {*} \alpha + \mathrm{pr} _ {2} ^ {*} \omega)} \end{array}
\]

从而 S 可积（Differentiable and Analytic Manifolds, R, 9.3.6）。

定理 5. 设 \(G\) 是一个李群芽，\(M\) 是一个 \(C^r (r\geqslant 2)\) 类流形，而 \(\alpha\) 是一个 \(\mathbf{C}^{r-1}\) 类 1 次微分形式，定义在 \(M\) 上并取值于 \(L(G)\)，满足 \(d\alpha +[\alpha ]^2 = 0\)。对所有 \(x\in M\) 和 \(g\in G\)，存在一个映射 \(f\)，它属于 \(\mathbf{C}^{r-1}\) 类，在 \(x\) 的某个开邻域上有定义，取值于 \(G\)，满足 \(f(x) = g\) 及 \(f^{-1}.df = \alpha\)。满足这些条件的两个映射在 \(x\) 的某个邻域上重合。

设 \(x \in M\) 且 \(g \in G\)。由引理 6（采用其记号）以及 Differentiable and Analytic Manifolds, R, 9.3.7，存在 M 中 \(x\) 的一个开邻域 U，以及一个 \(m \mapsto \phi(m) = (m, f(m))\) 的 \(C^{r-1}\) 类从 U 到 \(M \times G\) 的映射，使得 \(f(x) = g\) 且 \(\phi^*(\theta) = 0\)。于是

\[
\begin{array}{l l} f ^ {- 1}. d f = f ^ {*} (\omega) & \text {(§ 3, no. 18.9)} \\ = (\mathrm{pr} _ {2} \circ \phi) ^ {*} (\omega) & \text {(for \(f = \mathrm{pr} _ {2} \circ \phi\))} \\ = \phi^ {*} (\mathrm{pr} _ {1} ^ {*} \alpha - \theta) & \text {(Lemma 6)} \\ = \phi^ {*} (\mathrm{pr} _ {1} ^ {*} \alpha) & \text {(for \(\phi^ {*} (\theta) = 0\))} \\ = \alpha & \text {(for \(\mathrm{pr} _ {1} \circ \phi = \mathrm{Id} _ {\mathbf {U}}\))}. \end{array}
\]

设 \(f'\) 是从 U 到 G 的 \(\mathbf{C}^{r-1}\) 类映射，满足 \(f'(x) = g\) 和 \(f'^{-1}df' = \alpha\)。由 § 3, 18.9，\(ff'^{-1}\) 局部为常值，因而 \(f' = f\) 在 \(x\) 的某个邻域上成立。

命题 10. 设 M 是解析流形，g 是完备可赋范李代数，\(\alpha\) 是 M 上取值于 g 的解析 1 次微分形式，并具有以下性质：

\begin{enumerate}
\def\labelenumi{(\alph{enumi})}
\item
  对所有 \(m \in \mathbf{M}\)，\(\alpha_{m}\) 都是从 \(\mathrm{T}_{m}(\mathrm{M})\) 到 \(\mathfrak{g}\) 的同构；
\item
  \(d\alpha + [\alpha]^{2} = 0.\)
\end{enumerate}

那么，对所有 \(m_0 \in \mathbf{M}\)，存在一个开邻域 \(\mathbf{M}'\)，它是 \(m_0\) 在 \(\mathbf{M}\) 中的邻域，并且 \(\mathbf{M}'\) 上带有与 \(\mathbf{M}'\) 的流形结构相容的李群芽结构，单位元为 \(m_0\)，并具有以下性质：

\begin{enumerate}
\def\labelenumi{(\roman{enumi})}
\item
  \(\alpha_{m_{0}}\) 是从 L(M') 到 g 的同构；
\item
  微分形式 \(m \mapsto \alpha_{m_0}^{-1} \circ \alpha_m\) 是 \(M'\) 的左规范微分形式。若 \(M_1'\) 和 \(M_2'\) 是两个这样的群芽，则 \(M_1'\) 和 \(M_2'\) 有一个公共开子群芽。
\end{enumerate}

存在一个李群芽 G，使得 \(L(G) = g\)。令 \(m_{0} \in M\)。由定理 5，存在一个开邻域 \(M'\)（\(m_{0}\) 在 M 中），以及一个从 \(M'\) 到 G 的解析映射 f，使得 \(f(m_{0}) = e\) 且 \(f^{-1}.df = \alpha\)。于是 \(T_{m_{0}}(f) = \alpha_{m_{0}}\) 是从 \(T_{m_{0}}(M)\) 到 g 的同构；因此，缩小 \(M'\) 和 G 后，可以设 f 是从流形 \(M'\) 到流形 G 的同构。把 G 上的李群芽结构传递给 \(M'\)，所用的同构为 \(f^{-1}\)。于是 \(T_{m_{0}}(f)\) 成为从 \(L(M')\) 到 \(L(G) = g\) 的同构，从而得到 (i)。另一方面，若 \(\omega\) 表示 G 的左规范微分形式，则

\[
\begin{array}{r} \alpha_ {m _ {0}} ^ {- 1} \circ \alpha_ {m} = (\mathrm{T} _ {m _ {0}} f) ^ {- 1} \circ (f ^ {- 1}. d f) (m) \\ = (\mathrm{T} _ {m} f) ^ {- 1} \circ \omega (f (m)) \circ \mathrm{T} _ {m} f \end{array}
\]

从而 \(m \mapsto \alpha_{m_0}^{-1} \circ \alpha_m\) 是 \(\mathbf{M}'\) 的左规范微分形式。

令 \(\mathbf{M}^{\prime}\) 为 \(m_0\) 的一个开邻域，其上带有李群芽结构，单位元为 \(m_0\)，并具有与性质 (i) 和 (ii) 类似的性质。则 \(\alpha_{m_0}\) 既是从 \(\mathrm{L}(\mathbf{M}')\) 到 \(\mathfrak{g}\) 的同构，也是从 \(\mathrm{L}(\mathbf{M}'')\) 到 \(\mathfrak{g}\) 的同构，因而 \(\mathrm{L}(\mathbf{M}') = \mathrm{L}(\mathbf{M}'')\)。因此，缩小 \(\mathbf{M}'\) 和 \(\mathbf{M}''\) 后，可以设存在一个同构 \(\phi\)，它把群芽 \(\mathbf{M}'\) 映到群芽 \(\mathbf{M}''\)（no. 1, Corollary 1 to Theorem 1）。于是 \(\phi^{-1}.d\phi\) 是 \(\mathbf{M}'\) 的左规范微分。另一方面，令 \(\psi\) 为从流形 \(\mathbf{M}' \cap \mathbf{M}''\) 到李群芽 \(\mathbf{M}''\) 的规范嵌入；显然 \(\psi^{-1}.d\psi\) 是 \(\mathbf{M}''\) 的左规范微分的限制。因此有 \((\psi^{-1}.d\psi)(m) = \alpha_{m_0}^{-1} \circ \alpha_m = (\phi^{-1}.d\phi)(m)\)，对所有 \(m \in \mathbf{M}' \cap \mathbf{M}''\) 成立。所以 \(\phi\) 和 \(\psi\) 在 \(m_0\) 的某个邻域上重合（§ 3, 18.9）。这就证明了命题的最后一个断言。

推论。设 M 是有限维解析流形。~令 \(\omega_{1},\ldots,\omega_{n}\) 为 M 上取值为标量的解析 1 次微分形式，它们在 M 的每一点都线性无关，并且对所有 \(k=1,\ldots,n\)，\(\omega_{k}\) 都是 \(\omega_{i}\wedge\omega_{j}\) 的常系数线性组合。那么，对所有 \(m_{0}\in M\)，存在一个开邻域 \(M'\)，它是 \(m_{0}\) 在 M 中的邻域，并且带有定义在 \(M'\) 上、与 \(M'\) 的流形结构相容的李群芽结构，单位元为 \(m_{0}\)，使得 \(\omega_{1}|_{M'},\ldots,\omega_{n}|_{M'}\) 构成取值为标量的左不变 1 次微分形式空间在 \(M'\) 上的一组基。

若 \(\mathbf{M}_1'\) 和 \(\mathbf{M}_2'\) 是两个这样的群芽，则 \(\mathbf{M}_1'\) 和 \(\mathbf{M}_2'\) 有一个公共开子群芽。

令 \(\mathbf{X}_1, \ldots, \mathbf{X}_n\) 为 M 上的向量场，使得在 M 的每一点 \(m\)，\((\mathbf{X}_i)_m\) 构成 \(\mathrm{T}_m(\mathrm{M})\) 的一组与 \(((\omega_1)_m), \ldots, (\omega_n)_m)\) 对偶的基。这些向量场是解析的。由假设，存在 \(c_{ijk} \in \mathbf{K}\)（\(1 \leqslant i,j, k \leqslant n\)），使 \(c_{ijk} = -c_{jik}\) 且 \(d\omega_k = \sum_{i < j} c_{ijk}\omega_i \wedge \omega_j\)。由 Differentiable and Analytic Manifolds, R, 8.5.7，公式 (11)，

\[
\langle [ \mathrm{X} _ {i}, \mathrm{X} _ {j} ], \omega_ {k} \rangle = - (d \omega_ {k}) (\mathrm{X} _ {i}, \mathrm{X} _ {j}) = - \left(\sum_ {r <   s} c _ {r s k} \omega_ {r} \wedge \omega_ {s}\right) (\mathrm{X} _ {i}, \mathrm{X} _ {j}) = - c _ {i j k}
\]

从而 \([\mathrm{X}_i,\mathrm{X}_j] = -\sum_{k}c_{ijk}\mathrm{X}_k\)。于是 \(-c_{ijk}\) 是李代数 \(\mathfrak{g}\) 相对于基 \((e_1,\dots ,e_n)\) 的结构常数。对所有 \(m\in \mathbf{M}\)，令 \(\alpha_{m}\) 为从 \(\mathrm{T}_{m}(\mathrm{M})\) 到 \(\mathfrak{g}\) 的线性映射，它把 \((\mathrm{X}_1)_m\) 映到 \(e_1,\ldots ,(\mathrm{X}_n)_m\) 映到 \(e_n\)。于是 \(\alpha\) 是 \(\mathbf{M}\) 上的解析微分形式，次数为

1，取值于 \(\mathfrak{g}\)，且 \(\alpha_{m}\) 是从 \(\mathrm{T}_m(\mathbf{M})\) 到 \(\mathfrak{g}\) 的同构。另一方面，\(\alpha = \sum_{k=1}^{n} \omega_k e_k\)，因此

\[
d \alpha = \sum_ {k = 1} ^ {n} (d \omega_ {k}) e _ {k} = \sum_ {k = 1} ^ {n} \left(\sum_ {i <   j} c _ {i j k} \omega_ {i} \wedge \omega_ {j}\right) e _ {k}
\]

并且

\[
\begin{array}{r l} {[ \alpha ] ^ {2} = \sum_ {k = 1} ^ {n} [ \omega_ {k} e _ {k} ] ^ {2} + \sum_ {i <   j} (\omega_ {i} e _ {i}) \wedge (\omega_ {j} e _ {j})} & {\text {(§ 3, formula (30))}} \\ {= \sum_ {i <   j} (\omega_ {i} \wedge \omega _ {j}) [ e _ {i}, e _ {j} ]} \\ {= - \sum_ {k = 1} ^ {n} \sum_ {i <   j} (c _ {i j k} \omega_ {i} \wedge \omega_ {j}) e _ {k}} \\ {= - d \alpha.} \end{array}\tag{5}
\]

只需应用命题 10。

